% ===================== OBJETIVOS =====================
\section{Objetivos}

\subsection{Objetivo General}
Determinar la cantidad de bagazo de caña de azúcar colombiano necesaria para producir 100 toneladas por hora de vapor sobrecalentado a \SI{106}{barg} y \SI{545}{\celsius} en una caldera acuotubular con eficiencia del 94\%.

\subsection{Objetivos Específicos}
\begin{enumerate}
	\item Caracterizar el bagazo colombiano mediante su composición elemental y poder calorífico, utilizando datos experimentales de fuentes académicas nacionales.
	\item Realizar el balance de materia del ciclo agua-vapor considerando la purga continua del 2\%.
	\item Calcular el balance de energía térmica y determinar el consumo de bagazo.
	\item Expresar el resultado como el ratio de toneladas de vapor por tonelada de bagazo.
	\item Proporcionar una guía de simulación del proceso en Aspen Plus.
\end{enumerate}
